\documentclass[11pt]{article}
\usepackage[T1]{fontenc}
\usepackage{lmodern,amsmath,amssymb,amsthm,mathtools}
\usepackage[margin=1in]{geometry}
\usepackage[colorlinks=true,linkcolor=blue,urlcolor=blue,citecolor=blue]{hyperref}
\usepackage{microtype}
\newtheorem{theorem}{Theorem}
\newtheorem{lemma}{Lemma}
\newtheorem{corollary}{Corollary}
\newcommand{\N}{\mathbb N}
\newcommand{\R}{\mathbb R}
\title{Critical scaling of iterated rational\protect\\approximants to Euler's constant}
\author{Checked critical-limit note}
\date{12 October 2026}
\begin{document}
\maketitle
\begin{abstract}
We prove a joint limit for an established rational-approximation family when its degree and integer iteration count grow together. If $n\to\infty$ and $n^3 2^{-p}\to\lambda_0>0$, the original finite weighted ratio converges to an explicit harmonic-number series quotient. For $n=2^\ell$ and $p=3\ell$, elementary rational bounds show that the limit is strictly less than Euler's constant, so this sequence does not converge to that constant. Both statements, including their analytic dependencies and the original finite-sum definitions, are verified in Lean 4.32.2 with pinned Mathlib. The family and classical special-function ingredients have known sources. Historical novelty remains unresolved, and no irrationality result is claimed.
\end{abstract}

\section{Definitions and the completed result}
Let $H_0=0$, $H_m=\sum_{r=1}^m r^{-1}$, and
$\gamma=\lim_{m\to\infty}(H_m-\log m)$. For integers $n\ge1$, $p\ge0$, and $0\le k\le n$, set
\begin{align}
w_{n,p}(k)&=\binom nk^2\binom{n+k}{k}^{p}\frac1{k!},\label{eq:w}\\
a_{n,p}(k)&=(p+3)H_k-2H_{n-k}-pH_{n+k},\label{eq:a}\\
R_{n,p}&=\frac{\sum_{k=0}^n w_{n,p}(k)a_{n,p}(k)}{\sum_{k=0}^n w_{n,p}(k)}.\label{eq:R}
\end{align}
The weights are positive rational numbers. Thus the denominator is nonzero and $R_{n,p}$ is a rational number. For $x\ge0$, write
\[
S_0(x)=\sum_{j=0}^\infty\frac{x^j}{(j!)^2},\qquad
S_1(x)=\sum_{j=0}^\infty\frac{H_jx^j}{(j!)^2}.
\]
Both series converge; $S_0(x)\ge1$. Define, for $x>0$,
\begin{equation}
L(x)=3\gamma+\log x-2\frac{S_1(x)}{S_0(x)}.
\label{eq:L}
\end{equation}

\begin{theorem}[Joint critical limit]\label{thm:critical}
Let $n,p:\N\to\N$ satisfy $n_\ell\ge1$ for every $\ell$ and $n_\ell\to\infty$. If
\[
\lambda_\ell=\frac{n_\ell^3}{2^{p_\ell}}\longrightarrow\lambda_0>0,
\]
then the original weighted ratios in \eqref{eq:R} satisfy
\[
R_{n_\ell,p_\ell}\longrightarrow L(\lambda_0).
\]
\end{theorem}

\begin{theorem}[A concrete limit different from Euler's constant]\label{thm:concrete}
The sequence $R_{2^\ell,3\ell}$ converges to
\[
L(1)=3\gamma-2\frac{S_1(1)}{S_0(1)}<\gamma.
\]
In particular, it does not converge to $\gamma$.
\end{theorem}

These are the two completed results of this note. The quantifiers concern sequences of integer pairs; there is no real iteration count. No assumption about a denominator limit, parameter decay, summability, pointwise convergence or domination is added to Theorem~\ref{thm:critical}. Each is proved from its displayed hypotheses. Neither theorem addresses the rationality or irrationality of $\gamma$.

\section{The known family and prior work}
The family in \eqref{eq:w}--\eqref{eq:R} is the repeated-operator construction in Section~5 of Van Assche--Wolfs~\cite{vw}, evaluated at $x=1$. With $T_nf=(x^nf)^{(n)}/n!$, one has
\begin{align*}
T_n(x^k)&=\binom{n+k}{k}x^k,\\
T_n(x^k\log x)&=\binom{n+k}{k}x^k\bigl(\log x+H_{n+k}-H_k\bigr).
\end{align*}
Starting with the $p=0$ coefficients in their Proposition~2.2, iteration gives \eqref{eq:w}--\eqref{eq:a}. The construction itself is not claimed as new. The proved main asymptotic theorem in the checked arXiv version concerns $p=1$; the general-$p$ discussion in Section~5 proposes further asymptotic ideas, without a uniform remainder that establishes the present moving-$p$ limit.

The series in \eqref{eq:L} also have classical modified-Bessel interpretations~\cite{dlmfseries}, related to the Brent--McMillan computation of $\gamma$; see the survey~\cite{lagarias}. This note keeps the limit in explicit series form. It does not include a formal certificate for a Bessel identification or a strict gap for every $\lambda_0>0$.

The accompanying prior-art review compares theorem hypotheses and conclusions in the accessible versions of the cited family and related work. It identifies established construction, iteration and special-function ingredients. It does not establish historical novelty of the joint limit. Important gaps remain: the final 2025 journal text was not compared with its public preprint; the relevant 2026 Aptekarev survey was not available in full text; and wider growing-dimension hypergeometric literature has not been exhausted. Formal correctness does not resolve those historical questions. This public release is not a claim of external peer review or priority.

\section{Endpoint weights and summable bounds}
For $n\ge1$, define the zero-extended endpoint weights
\[
r_{n,p}(j)=
\begin{cases}
w_{n,p}(n-j)/w_{n,p}(n),&0\le j\le n,\\
0,&j>n.
\end{cases}
\]
Put $\lambda=n^3/2^p$ and $v_x(j)=x^j/(j!)^2$. Factorial cancellation gives the exact identity
\begin{equation}
r_{n,p}(j)=v_\lambda(j)\prod_{r=0}^{j-1}
\left(1-\frac rn\right)^3
\left(\frac{1-r/n}{1-r/(2n)}\right)^p
\qquad(0\le j\le n).
\label{eq:product}
\end{equation}
Every factor in this range lies in $[0,1]$. Therefore, including $j>n$,
\begin{equation}
0\le r_{n,p}(j)\le v_\lambda(j).
\label{eq:domination}
\end{equation}
The zero-extension is essential: it excludes spurious terms that truncated natural subtraction could otherwise introduce.

Under the hypotheses of Theorem~\ref{thm:critical}, the exact logarithmic identity
\[
\frac{p_\ell}{n_\ell}=
\frac{3\log(n_\ell)/n_\ell-\log(\lambda_\ell)/n_\ell}{\log2}
\]
shows $p_\ell/n_\ell\to0$, and hence $(p_\ell+3)/n_\ell\to0$.
For fixed $r$ and eventually $r<n_\ell$, the factor $c_{\ell,r}$ in \eqref{eq:product} satisfies
\[
0\le1-c_{\ell,r}\le\frac{3r}{n_\ell}
+\frac{p_\ell r}{2n_\ell-r}
\le\frac{(p_\ell+3)r}{n_\ell}.
\]
Here we use $1-u^m\le m(1-u)$ for $0\le u\le1$ and integer $m\ge0$, and $1-ab\le(1-a)+(1-b)$ for $a,b\in[0,1]$. Each fixed factor tends to one. Since $n_\ell\to\infty$, every fixed $j$ eventually lies below the endpoint. Equation~\eqref{eq:product} then proves
\begin{equation}
r_{n_\ell,p_\ell}(j)\longrightarrow v_{\lambda_0}(j)
\qquad\text{for every fixed }j.
\label{eq:pointwise}
\end{equation}

Set $M=\lambda_0+1$. Eventually $0\le\lambda_\ell\le M$, with one index threshold valid for all $j$. The three series
\[
\sum_jv_M(j),\qquad \sum_jH_jv_M(j),\qquad \sum_j j^2v_M(j)
\]
converge. For example $H_j\le j$, polynomially weighted factorial series converge, and $j^2v_M(j)=M v_M(j-1)$ for $j\ge1$. Dominated convergence applied to \eqref{eq:domination}--\eqref{eq:pointwise} gives
\begin{align}
D_\ell:=\sum_jr_{n_\ell,p_\ell}(j)&\longrightarrow S_0(\lambda_0),\label{eq:D}\\
\sum_jH_jr_{n_\ell,p_\ell}(j)&\longrightarrow S_1(\lambda_0).\label{eq:H}
\end{align}
The bounds are uniform in the summation index after a single sufficiently large sequence index.

\section{The numerator and the critical limit}
Write
\[
C_{n,p}=3H_n-p(H_{2n}-H_n).
\]
For $0\le j\le n$, the original coefficient has the exact decomposition
\begin{align}
a_{n,p}(n-j)&=C_{n,p}-2H_j+\delta_{n,p}(j),\label{eq:decomp}\\
\delta_{n,p}(j)&=-(p+3)(H_n-H_{n-j})
+p(H_{2n}-H_{2n-j}).\notag
\end{align}
We need a bound valid even at $j=n$. For $0\le r<n$,
\[
\frac1{n-r}\le\frac{r+1}{n},
\]
since $n\le(r+1)(n-r)$. Summing yields
\[
0\le H_n-H_{n-j}\le\frac{j(j+1)}{2n}\le\frac{j^2}{n}
\qquad(1\le j\le n),
\]
and the zero case is immediate. Also
$0\le H_{2n}-H_{2n-j}\le H_n-H_{n-j}$. Thus
\begin{equation}
|\delta_{n,p}(j)|\le\frac{(p+3)j^2}{n}.
\label{eq:correction}
\end{equation}
Extend $r_{n,p}(j)\delta_{n,p}(j)$ by zero for $j>n$. Equations~\eqref{eq:domination} and \eqref{eq:correction} give
\[
|r_{n,p}(j)\delta_{n,p}(j)|
\le\frac{p+3}{n}\,j^2v_\lambda(j).
\]
For every fixed $j$, this tends to zero along the critical sequence. Eventually $(p_\ell+3)/n_\ell\le1$ and $\lambda_\ell\le M$, so $j^2v_M(j)$ is a summable majorant independent of $\ell$. Dominated convergence proves
\begin{equation}
\sum_j r_{n_\ell,p_\ell}(j)\delta_{n_\ell,p_\ell}(j)\longrightarrow0.
\label{eq:E}
\end{equation}

The classical harmonic remainder $h_n=H_n-\log n-\gamma$ satisfies $0\le h_n\le1/n$ for $n\ge1$. Consequently
\[
C_{n,p}=3\gamma+\log\lambda+3h_n+p h_n-p h_{2n}.
\]
Because $n_\ell\to\infty$ and $p_\ell/n_\ell\to0$, it follows that
\begin{equation}
C_{n_\ell,p_\ell}\longrightarrow3\gamma+\log\lambda_0.
\label{eq:C}
\end{equation}
Reversing the finite sum and dividing numerator and denominator by the positive endpoint weight gives exactly
\[
R_{n,p}=C_{n,p}
-2\frac{\sum_jH_jr_{n,p}(j)}{\sum_jr_{n,p}(j)}
+\frac{\sum_jr_{n,p}(j)\delta_{n,p}(j)}{\sum_jr_{n,p}(j)}.
\]
The limit denominator is $S_0(\lambda_0)\ge1$. Combining \eqref{eq:D}, \eqref{eq:H}, \eqref{eq:E} and \eqref{eq:C} proves Theorem~\ref{thm:critical}.

\section{An explicit strict gap}
For $n_\ell=2^\ell$ and $p_\ell=3\ell$, one has $n_\ell\to\infty$ and $n_\ell^3/2^{p_\ell}=1$. Theorem~\ref{thm:critical} therefore gives convergence to $L(1)$. We prove $L(1)<\gamma$ without using Bessel functions.

The first two terms of the series
\[
\log\frac{1+x}{1-x}=2\sum_{m=0}^\infty\frac{x^{2m+1}}{2m+1}
\qquad(0<x<1)
\]
at $x=1/3$ give $\log2\ge56/81$. The strict harmonic upper bound for $\gamma$ yields
\[
\gamma<H_{16}-\log16
\le\frac{2436559}{720720}-\frac{224}{81}
=\frac{3991111}{6486480}<\frac58.
\]
On the other hand, exact rational arithmetic gives
\[
\sum_{j=0}^{3}\frac{H_j-5/8}{(j!)^2}=\frac1{432}>0.
\]
Every subsequent summand is nonnegative, since $H_j\ge1$ for $j\ge1$. Summability and $S_0(1)>0$ imply
\[
\frac{S_1(1)}{S_0(1)}>\frac58>\gamma.
\]
Hence $L(1)=3\gamma-2S_1(1)/S_0(1)<\gamma$. Uniqueness of limits shows that $R_{2^\ell,3\ell}$ cannot also converge to $\gamma$, proving Theorem~\ref{thm:concrete}.

\section{Formal scope and reproducibility}
The accompanying twelve local Lean modules form the complete local source closure, including the explicit statement fixture and concrete strict-gap proof. The general theorem, in the namespace \texttt{EulerScalingAnalytic}, is
\begin{center}
\texttt{tendsto\_approximant\_critical}.
\end{center}
The file \texttt{ExpectedSpec.lean} states its finite binomial weights, harmonic coefficients, two finite sums, integer-pair scaling hypothesis and limit series explicitly, and proves agreement by unfolding the exported definitions. The concrete theorem file proves convergence to $L(1)$, $L(1)<\gamma$, and failure to converge to $\gamma$.

All analytic dependencies and exact finite reindexings used above are proved in this closure. A fresh-output rebuild of all twelve modules passed, reusing pinned upstream caches. The final axiom audits list only \texttt{propext}, \texttt{Classical.choice} and \texttt{Quot.sound}. The sources contain no admitted proof, custom axiom, unsafe declaration, native-evaluation shortcut or implementation override. No source rebuild of all of Mathlib or separate-kernel replay is claimed.

The checked versions are Lean 4.32.2, commit \texttt{f3b06c7}, and Mathlib v4.32.2, commit \texttt{905b958}. Full revisions, the dependency manifest, build script, source checksums and compiler output accompany the unchanged checked Lean sources.

Quantitative finite-parameter estimates, the diverging-parameter regime, the Bessel identity and a general strict gap for all positive $\lambda_0$ are outside this certificate and are not additional completed results of this note. No numerical experiment is needed. The certificate implies no irrationality, arithmetic-quality or historical-novelty conclusion.

\begin{thebibliography}{9}
\bibitem{vw} W. Van Assche and T. Wolfs, \emph{Rational approximation of Euler's constant using multiple orthogonal polynomials}, Combinatorics and Number Theory 14 (2025), 141--162. \href{https://doi.org/10.2140/cnt.2025.14.141}{doi:10.2140/cnt.2025.14.141}. Checked version: \href{https://arxiv.org/html/2404.09799v3}{arXiv:2404.09799v3}.
\bibitem{lagarias} J. C. Lagarias, \emph{Euler's constant: Euler's work and modern developments}, Bull. Amer. Math. Soc. 50 (2013), 527--628. \href{https://arxiv.org/abs/1303.1856}{arXiv:1303.1856}.
\bibitem{dlmfseries} NIST Digital Library of Mathematical Functions, \href{https://dlmf.nist.gov/10.31#E2}{Equation 10.31.2} and \href{https://dlmf.nist.gov/10.25#E2}{Equation 10.25.2}.
\bibitem{aptekarev} A. I. Aptekarev, \emph{Multiple orthogonal polynomials and strong Szeg\H{o} asymptotics}, Russian Mathematical Surveys 81(2) (2026), 209--273. \href{https://www.mathnet.ru/eng/rm10281}{doi:10.4213/rm10281e}. Bibliographic record checked; full text not reviewed.
\end{thebibliography}
\end{document}
